\documentclass[a4paper,10pt]{letter}

\usepackage[T1]{fontenc}
\usepackage[utf8]{inputenc}
\usepackage{lmodern}
\usepackage[hidelinks]{hyperref}
\usepackage{microtype}
\usepackage[margin=1in]{geometry}

\signature{Omar Iskandarani\\
Independent Researcher, Groningen\\
The Netherlands\\
ORCID: 0009-0006-1686-3961\\
Email: \href{mailto:info@omariskandarani.com}{info@omariskandarani.com}}

\address{Omar Iskandarani\\
Vinkenstraat 86A\\
9713 TK Groningen\\
The Netherlands}

\date{\today}

\begin{document}

    \begin{letter}{Editors\\
    \textit{International Journal of Theoretical Physics}\\
    Springer Nature}

\opening{Dear Editors,}

Please consider the enclosed manuscript entitled \textbf{“Phase Locking and Discrete Mode Selection in Closed Circulation Loops”} for publication in the \textit{International Journal of Theoretical Physics}.

The manuscript presents a rigorous analytical study of discrete mode selection in a minimal delayed phase oscillator representing a closed circulation loop with finite propagation time. We show that integer-indexed spectral branches emerge directly from the stability structure of a differenced delay differential equation, without the need to impose spatial boundary conditions or quantization assumptions.

The principal results include:

\begin{itemize}
    \item Derivation of the exact characteristic equation governing linearized phase perturbations.
    \item Proof of unconditional stability for one asymptotic family of branches and instability of intermediate branches in the strong-feedback regime.
    \item Establishment of an asymptotic frequency spacing law $\Delta\Omega \approx 2\pi/\tau$ and the corresponding spectral density scaling $\rho(\Omega) \approx \tau/(2\pi)$.
\end{itemize}

These findings clarify how discrete spectral families can arise purely from temporal nonlocality and delay-induced stability filtering in continuous dynamical systems. The analysis is fully classical and confined to nonlinear dynamics and delay theory; no modification of established physical frameworks is proposed. Rather, the work isolates a structural mechanism within delayed circulation systems that may be relevant in theoretical contexts where finite propagation time plays a central role.

We believe this contribution aligns with the journal’s emphasis on mathematically substantive results that reveal structural connections across areas of theoretical physics.

Thank you for your consideration.

\closing{Sincerely,}

    \end{letter}

\end{document}